\documentclass{amsart}
% For dealing with references we use the comment environment
\usepackage{verbatim}
\newenvironment{reference}{\comment}{\endcomment}
%\newenvironment{reference}{}{}
\newenvironment{slogan}{\comment}{\endcomment}
\newenvironment{history}{\comment}{\endcomment}

% For commutative diagrams we use Xy-pic
\usepackage[all]{xy}

% We use 2cell for 2-commutative diagrams.
\xyoption{2cell}
\UseAllTwocells

% We use multicol for the list of chapters between chapters
\usepackage{multicol}

% This is generally recommended for better output
\usepackage{lmodern}
\usepackage[T1]{fontenc}

% For cross-file-references

% Package for hypertext links:
\usepackage{hyperref}

% For any local file, say "hello.tex" you want to link to please

% Theorem environments.
%
\theoremstyle{plain}
\newtheorem{theorem}[subsection]{Theorem}
\newtheorem{proposition}[subsection]{Proposition}
\newtheorem{lemma}[subsection]{Lemma}

\theoremstyle{definition}
\newtheorem{definition}[subsection]{Definition}
\newtheorem{example}[subsection]{Example}
\newtheorem{exercise}[subsection]{Exercise}
\newtheorem{situation}[subsection]{Situation}

\theoremstyle{remark}
\newtheorem{remark}[subsection]{Remark}
\newtheorem{remarks}[subsection]{Remarks}

\numberwithin{equation}{subsection}

% Macros
%
\def\lim{\mathop{\mathrm{lim}}\nolimits}
\def\colim{\mathop{\mathrm{colim}}\nolimits}
\def\Spec{\mathop{\mathrm{Spec}}}
\def\Hom{\mathop{\mathrm{Hom}}\nolimits}
\def\Ext{\mathop{\mathrm{Ext}}\nolimits}
\def\SheafHom{\mathop{\mathcal{H}\!\mathit{om}}\nolimits}
\def\SheafExt{\mathop{\mathcal{E}\!\mathit{xt}}\nolimits}
\def\Sch{\mathit{Sch}}
\def\Mor{\mathop{\mathrm{Mor}}\nolimits}
\def\Ob{\mathop{\mathrm{Ob}}\nolimits}
\def\Sh{\mathop{\mathit{Sh}}\nolimits}
\def\NL{\mathop{N\!L}\nolimits}
\def\CH{\mathop{\mathrm{CH}}\nolimits}
\def\proetale{{pro\text{-}\acute{e}tale}}
\def\etale{{\acute{e}tale}}
\def\QCoh{\mathit{QCoh}}
\def\Ker{\mathop{\mathrm{Ker}}}
\def\Im{\mathop{\mathrm{Im}}}
\def\Coker{\mathop{\mathrm{Coker}}}
\def\Coim{\mathop{\mathrm{Coim}}}

% Boxtimes
%
\DeclareMathSymbol{\boxtimes}{\mathbin}{AMSa}{"02}

%
% Macros for moduli stacks/spaces
%
\def\QCohstack{\mathcal{QC}\!\mathit{oh}}
\def\Cohstack{\mathcal{C}\!\mathit{oh}}
\def\Spacesstack{\mathcal{S}\!\mathit{paces}}
\def\Quotfunctor{\mathrm{Quot}}
\def\Hilbfunctor{\mathrm{Hilb}}
\def\Curvesstack{\mathcal{C}\!\mathit{urves}}
\def\Polarizedstack{\mathcal{P}\!\mathit{olarized}}
\def\Complexesstack{\mathcal{C}\!\mathit{omplexes}}
% \Pic is the operator that assigns to X its picard group, usage \Pic(X)
% \Picardstack_{X/B} denotes the Picard stack of X over B
% \Picardfunctor_{X/B} denotes the Picard functor of X over B
\def\Pic{\mathop{\mathrm{Pic}}\nolimits}
\def\Picardstack{\mathcal{P}\!\mathit{ic}}
\def\Picardfunctor{\mathrm{Pic}}
\def\Deformationcategory{\mathcal{D}\!\mathit{ef}}

\setlength{\emergencystretch}{3em}
\begin{document}
\title[Stacks Project, Tag 09Z2]{275 --- Cohomological descent and its spectral sequence for surjective finite scheme morphisms}
\maketitle
\noindent Copyright (C) 2005--2025 Johan de Jong. Stacks Project authors. Source: \href{https://stacks.math.columbia.edu/tag/09Z2}{Tag 09Z2}.

\noindent Editorial difficulty note (not author text). Difficulty H2: The proof constructs the augmented cosimplicial sheaf complex attached to the entire Cech nerve of a finite surjection. Exactness must be proved on geometric stalks by identifying the nerve with powers of a finite nonempty fibre, then finite pushforward acyclicity must be used to obtain the asserted descent spectral sequence..

\noindent Editorial provenance note (not author text). History: excerpt from revision \texttt{a0444\allowbreak 6e57e\allowbreak c1fbc\allowbreak 252a8\allowbreak 71afc\allowbreak ec775\allowbreak 2fb28\allowbreak 07b14}, etale-cohomology.tex, lines 8411--8485. Reference presentation was adapted; original-author context and core support blocks were retained or added. The mathematical wording is unchanged. Licensed under GNU FDL 1.2 or later; no invariant sections or cover texts. See ../licenses/COPYING. Context blocks reproduce author definitions and supporting results; they are not additional counted problems.

\begin{lemma}
\label{lemma-cohomological-descent-finite}
Let $f : X \to Y$ be a surjective finite morphism of schemes.
Set $f_n : X_n \to Y$ equal to the $(n + 1)$-fold fibre product
of $X$ over $Y$. For $\mathcal{F} \in \textit{Ab}(Y_\etale)$ set
$\mathcal{F}_n = f_{n, *}f_n^{-1}\mathcal{F}$. There is an exact
sequence
$$
0 \to \mathcal{F} \to \mathcal{F}_0 \to \mathcal{F}_1 \to
\mathcal{F}_2 \to \ldots
$$
on $Y_\etale$. Moreover, there is a spectral sequence
$$
E_1^{p, q} = H^q_\etale(X_p, f_p^{-1}\mathcal{F})
$$
converging to $H^{p + q}(Y_\etale, \mathcal{F})$.
This spectral sequence is functorial in $\mathcal{F}$.
\end{lemma}

\begin{proof}
If we prove the first statement of the lemma, then we obtain a spectral
sequence with $E_1^{p, q} = H^q_\etale(Y, \mathcal{F}_p)$ converging
to $H^{p + q}(Y_\etale, \mathcal{F})$, see
Derived Categories, Lemma \href{https://stacks.math.columbia.edu/tag/015J}{lemma two ss complex functor (Tag 015J)}.
On the other hand, since
$R^if_{p, *}f_p^{-1}\mathcal{F} = 0$ for $i > 0$
(Proposition \href{https://stacks.math.columbia.edu/tag/03QP}{proposition finite higher direct image zero (Tag 03QP)})
we get
$$
H^q_\etale(X_p, f_p^{-1}\mathcal{F}) =
H^q_\etale(Y, f_{p, *}f_p^{-1} \mathcal{F}) =
H^q_\etale(Y, \mathcal{F}_p)
$$
by Proposition \href{https://stacks.math.columbia.edu/tag/03QC}{proposition leray (Tag 03QC)}
and we get the spectral sequence of the lemma.

\medskip\noindent
To prove the first statement of the lemma, observe that
$X_n$ forms a simplicial scheme over $Y$, see
Simplicial, Example \href{https://stacks.math.columbia.edu/tag/016E}{example fibre products simplicial object (Tag 016E)}.
Observe moreover, that for each of the projections
$d_j : X_{n + 1} \to X_n$ there is a map
$d_j^{-1} f_n^{-1}\mathcal{F} \to f_{n + 1}^{-1}\mathcal{F}$.
These maps induce maps
$$
\delta_j : \mathcal{F}_n \to \mathcal{F}_{n + 1}
$$
for $j = 0, \ldots, n + 1$. We use the alternating sum of these maps
to define the differentials $\mathcal{F}_n \to \mathcal{F}_{n + 1}$.
Similarly, there is a canonical augmentation $\mathcal{F} \to \mathcal{F}_0$,
namely this is just the canonical map $\mathcal{F} \to f_*f^{-1}\mathcal{F}$.
To check that this sequence of sheaves is an exact complex it suffices
to check on stalks at geometric points (Theorem \href{https://stacks.math.columbia.edu/tag/03PU}{theorem exactness stalks (Tag 03PU)}).
Thus we let $\overline{y} : \Spec(k) \to Y$ be a geometric point. Let
$E = \{\overline{x} : \Spec(k) \to X \mid f(\overline{x}) = \overline{y}\}$.
Then $E$ is a finite nonempty set and we see that
$$
(\mathcal{F}_n)_{\overline{y}} =
\bigoplus\nolimits_{e \in E^{n + 1}} \mathcal{F}_{\overline{y}}
$$
by Proposition \href{https://stacks.math.columbia.edu/tag/03QP}{proposition finite higher direct image zero (Tag 03QP)}
and Lemma \href{https://stacks.math.columbia.edu/tag/03Q1}{lemma stalk pullback (Tag 03Q1)}.
Thus we have to see that given an abelian group $M$ the sequence
$$
0 \to M \to \bigoplus\nolimits_{e \in E} M \to
\bigoplus\nolimits_{e \in E^2} M \to \ldots
$$
is exact. Here the first map is the diagonal map and the map
$\bigoplus_{e \in E^{n + 1}} M  \to \bigoplus_{e \in E^{n + 2}} M$
is the alternating sum of the maps induced by the $(n + 2)$
projections $E^{n + 2} \to E^{n + 1}$. This can be shown directly
or deduced by applying Simplicial, Lemma
\href{https://stacks.math.columbia.edu/tag/019P}{lemma fibre products simplicial object w section (Tag 019P)}
to the map $E \to \{*\}$.
\end{proof}
\end{document}
